\section{Safeguarding the Mysteries}

Gornahoor has often emphasized that spiritual questions, or debates over the superiority this of that tradition, can be resolved neither by personal predilections nor by empirical and historical considerations. These issues can only be addressed from the understanding of metaphysical principles. Gornahoor finds it curious, though symptomatic of the contemporary human situation, that everyone feels competent to opine on spiritual and political matters, although they would not dare to say anything at all about quantum physics or algebraic topology, topics much simpler to understand.

One such principle is that knowledge is being, ``to know is to be''. In order to know something of spiritual depth, one must become deep oneself. Schuon wrote a book titled The Transcendent Unity of Religions; although on the human plane, different religions may diverge widely, at the level of principle, they must needs converge. Thus, for example, the Catholic monk Thomas Merton can rightly claim that he has more in common with a Zen monk than with the average Catholic in the pew.

Hence, Gornahoor has always emphasized the importance of spiritual practice. This is what separate the Metaphysician and the Hermetist from the ordinary philosopher. The great Hermetist of the 20th century, Valentin Tomberg, in the Meditations on the 22 Major Arcana of the Tarot makes this clear in the following passage: (p 122, my translation)

\begin{quotationx}
The goal of spiritual exercises is \emph{depth}. It is necessary to become deep in order to be able to achieve the experience and knowledge of deep things. And it is symbolism which is the language of depth --- so that these are the arcana expressed by the symbols that are the means and the goal of spiritual exercises, from which the living Tradition of Hermetic Philosophy is composed.

Common spiritual exercises create the common link which unites Hermetists. It is not common knowledge which unites them but rather spiritual exercises and the experiences they entail. If three people from different countries who had used \textbf{Moses's Genesis}, \textbf{John's Gospel}, and \textbf{Ezekiel's Vision} as the subjects of spiritual exercises for several years were to meet, they would do it as brothers even though one knew the history of humanity, the other had the science of healing, and the third was a deep cabbalist. What they \emph{know} is the result of \emph{personal} experience and direction, whereas the \emph{depth}, the \emph{level} that they have reached --- without regard to the aspect and extent of the knowledge that has been gained --- is what they have in \emph{common}. Hermetism, the Hermetic Tradition, is first of all and most of all, a certain degree of depth, a certain level of consciousness. And it is spiritual exercises that safeguard it.
\end{quotationx}



\flrightit{Posted on 2010-08-07 by Cologero}
